% -----------------------------------------------------------------------
% System Design
% -----------------------------------------------------------------------

\subsection{Overview}
\label{sec:design:overview}

VisionServe is a single statically-linked Go binary that exposes a
language-agnostic HTTP REST API for computer vision inference.
Clients send a \texttt{POST /api/predict} request carrying an image (base64
or multipart) together with a task identifier, a model name, and an optional
prompt; the server returns a JSON object conforming to the unified
\texttt{Result} schema described in Section~\ref{sec:design:result}.
The same endpoint handles all six supported task types --- detection,
segmentation, open-vocabulary detection, classification, depth estimation, and
embedding --- so no client-side multiplexing is required.

Internally the server is composed of four coordinating subsystems: a
\emph{model registry} that resolves model names and validates manifests at
startup, a \emph{lifecycle manager} that owns all ONNX Runtime sessions, an
\emph{inference engine} that drives the pre-process/infer/post-process
pipeline, and a \emph{router/handler} layer that parses requests and serialises
responses.
All heavy resources (GPU memory, session handles) are allocated during startup
and reused across requests; no resource allocation occurs in the request hot
path.

\begin{figure}[t]
  \centering
  \begin{tikzpicture}[
      font=\small,
      node distance=0.55cm and 0.9cm,
      box/.style={draw, rounded corners=3pt, minimum width=2.2cm,
                  minimum height=0.65cm, align=center, fill=blue!6},
      ebox/.style={draw, rounded corners=3pt, minimum width=2.2cm,
                   minimum height=0.65cm, align=center, fill=orange!10},
      mbox/.style={draw, rounded corners=3pt, minimum width=2.0cm,
                   minimum height=0.58cm, align=center, fill=green!8,
                   font=\footnotesize},
      groupbox/.style={draw=gray!60, dashed, rounded corners=4pt,
                       inner sep=6pt, fill=gray!4},
      arr/.style={-{Stealth[length=4pt]}, thick},
  ]

  %% HTTP client
  \node[box] (client) {HTTP Client};

  %% Server internals group
  \node[box, right=1.0cm of client] (router)  {Router / Handler};
  \node[box, right=0.8cm of router] (engine)  {Engine};
  \node[box, above=0.5cm of engine] (registry){Registry};
  \node[box, below=0.5cm of engine] (lifecycle){Lifecycle\\Manager};

  \begin{scope}[on background layer]
    \node[groupbox, fit=(router)(registry)(lifecycle)(engine),
          label={[font=\footnotesize,gray]above:VisionServe Server}] (server) {};
  \end{scope}

  %% ORT group
  \node[ebox, right=1.0cm of engine] (ort) {ONNX Runtime};
  \node[ebox, above=0.35cm of ort,   font=\footnotesize] (ep1) {TensorRT / CUDA};
  \node[ebox, below=0.35cm of ort,   font=\footnotesize] (ep2) {CoreML / DirectML\\OpenVINO / CPU};

  \begin{scope}[on background layer]
    \node[groupbox, fit=(ort)(ep1)(ep2),
          label={[font=\footnotesize,gray]above:Accelerators}] (ortgroup) {};
  \end{scope}

  %% Model interface group (below Engine)
  \node[mbox, below=1.2cm of lifecycle] (mplain) {\texttt{Model}\\(pre/post)};
  \node[mbox, right=0.6cm of mplain]   (mpipe)  {\texttt{PipelineModel}\\(chained)};

  \begin{scope}[on background layer]
    \node[groupbox, fit=(mplain)(mpipe),
          label={[font=\footnotesize,gray]below:Model Packages}] (mgroup) {};
  \end{scope}

  %% Arrows
  \draw[arr] (client)   -- node[above,font=\footnotesize]{POST} (router);
  \draw[arr] (router)   -- (engine);
  \draw[arr] (registry) -- (engine);
  \draw[arr] (lifecycle)-- (engine);
  \draw[arr] (engine)   -- (ort);
  \draw[arr,gray!70] (ort) -- (ep1);
  \draw[arr,gray!70] (ort) -- (ep2);
  \draw[arr,dashed]  (engine)    -- (lifecycle);
  \draw[arr,dashed]  (lifecycle) -- (mplain);
  \draw[arr,dashed]  (lifecycle) -- (mpipe);
  \draw[arr] (router) -- node[right,font=\footnotesize]{Result} ++(0,-2.8) -| (client);

  \end{tikzpicture}
  \caption{High-level architecture of VisionServe.
           Solid arrows denote data flow; dashed arrows denote ownership or
           interface delegation.
           The Lifecycle Manager owns all ONNX sessions and is the sole
           gateway between model logic and the ONNX Runtime.}
  \label{fig:architecture}
\end{figure}

\subsection{Model Registry and Manifest}
\label{sec:design:registry}

Every model is described by a YAML manifest that declares all metadata
necessary for startup validation and session creation.
The mandatory fields are: \texttt{name} (unique identifier),
\texttt{task} (one of the six supported task types),
\texttt{license} (SPDX identifier),
\texttt{execution\_providers} (ordered list of preferred EPs),
preprocessing parameters (input resolution, normalisation statistics, and
letterbox strategy), and either a single \texttt{model\_file} path for
single-session models or a \texttt{files} map from role names to file paths
for multi-session models.
Roles are string keys agreed upon between the manifest and the model package:
for example, MobileSAM declares \texttt{encoder} and \texttt{decoder} roles;
Grounded-SAM declares \texttt{det}, \texttt{encoder}, and \texttt{decoder}.

At startup the registry parser reads every manifest in the model directory and
enforces the license allowlist: only \texttt{Apache-2.0}, \texttt{MIT}, and
BSD variants are accepted.
Any manifest bearing a non-allowlisted license --- including AGPL-3.0, as
carried by Ultralytics YOLO~\cite{ultralytics2023yolo,agpllicense} --- is
rejected with a descriptive error before any session is created.
This approach provides a strong audit boundary: downstream adopters can trust
that every model loaded by VisionServe carries a permissive license, regardless
of where the ONNX file was sourced.

\subsection{Lifecycle Management}
\label{sec:design:lifecycle}

ONNX Runtime sessions are expensive GPU resources: each session allocates
device memory proportional to model weight size and holds an execution context
that spans the full session lifetime.
Creating sessions on demand inside request handlers would cause repeated
allocation/deallocation cycles, unpredictable latency spikes, and potential
VRAM fragmentation.

\texttt{lifecycle.Manager} solves this by owning the full lifecycle of every
session.
At startup the manager creates one session per role per model, holds references
for the duration of the process, and releases them only on server shutdown.
Concurrent HTTP requests access sessions through a mutex-protected pool: a
request acquires a session handle, performs inference, and returns it to the
pool, ensuring that a session is never used by two goroutines simultaneously.
Neither model packages nor request handlers hold any reference to an ORT
session object; all access is mediated through the manager's typed API.
This single ownership invariant eliminates the most common class of
GPU resource leak in serving systems and is straightforward to audit.

\subsection{Model Interfaces}
\label{sec:design:interfaces}

VisionServe defines two Go interfaces that cover the full range of CV models.

\textbf{Model} is the interface for plain single-session models.
An implementation provides \texttt{Preprocess(img) (tensor, PreprocessMeta,
error)} and \texttt{Postprocess(output, meta) (Result, error)}.
The inference engine handles everything in between: it acquires a session
from the lifecycle manager, feeds the preprocessed tensor, collects the raw
output tensors, and calls \texttt{Postprocess}.
Model packages contain no inference code; they contain only numerical
transformations on image data and model output arrays.
This split keeps model packages testable in isolation --- unit tests exercise
pre/postprocess without an ONNX session --- and keeps the hot path in the
engine, where batching and concurrency controls can be applied uniformly.

\textbf{PipelineModel} is the interface for models that require a
\texttt{Prompt} (bounding box, point set, or text string) or that orchestrate
multiple ONNX sessions in sequence.
An implementation provides \texttt{Roles() []string}, which declares the role
keys the model expects in the session pool, and \texttt{Infer(img, prompt,
Runner) (Result, error)}, which drives the full pipeline.
The \texttt{Runner} gateway exposes a single method \texttt{Run(role, tensor)
([]Tensor, error)} that dispatches to the lifecycle manager by role name;
the pipeline model calls \texttt{Runner.Run} once per session in its chain and
assembles the final result.
Critically, the pipeline model never creates or frees sessions; it merely
names them.
This design supports arbitrarily deep chains --- GroundingDINO boxes flowing
into the MobileSAM encoder then decoder in Grounded-SAM --- without any
special-casing in the core engine.

\subsection{Unified Result Schema}
\label{sec:design:result}

All six task types return a single \texttt{Result} struct:
\begin{center}
\small
\texttt{Result\{task, model, Detections, Masks, Classifications,}\\
\texttt{DepthMap, Embeddings, DurationMs\}}
\end{center}
Fields irrelevant to a given task are null or empty.
\texttt{Detection} carries a \texttt{BBox} in the format \([x, y, w, h]\) in
\emph{original image coordinates}: the postprocess step uses the
\texttt{PreprocessMeta} produced during preprocessing --- which records the
original image dimensions and any letterbox padding offsets --- to map model-
space coordinates back to pixel-space before populating the result.
This mapping is done once, inside \texttt{Postprocess}, so no consumer of
\texttt{Result} ever needs to account for scaling or padding.

Masks are encoded as column-major (Fortran-order) run-length encoding (RLE) in
the COCO uncompressed style.
Column-major order matches the internal memory layout of SAM decoder outputs
and avoids a transpose when encoding, keeping the mask serialisation path
allocation-free.
A single schema simplifies client code: a generic display widget or
post-processing script works across all tasks without branching on task type.

\subsection{Execution Provider Fallback}
\label{sec:design:ep}

Hardware diversity is a defining constraint of edge CV deployment: the same
model may run on a Jetson module (TensorRT + CUDA), a developer laptop (CUDA
or CoreML), an Intel NUC (OpenVINO), or a Raspberry Pi (CPU only).
Rather than shipping separate binaries, VisionServe resolves this at session
creation time through a manifest-declared EP priority list such as
\texttt{[tensorrt, cuda, cpu]}.
The engine iterates the list in order; if ONNX Runtime reports that a given EP
is unavailable in the current ORT build, it is silently skipped and the next
is tried.
CPU is always appended as the final fallback regardless of what the manifest
declares, so inference always succeeds on any host.
The allowlisted EPs are: \texttt{tensorrt}, \texttt{cuda}, \texttt{coreml},
\texttt{directml}, \texttt{openvino}, and \texttt{cpu}; any EP name not in
this set is rejected at manifest parse time to prevent misconfiguration.
The fallback chain is fully transparent: the \texttt{Result} struct records
which EP was ultimately selected, so operators can verify hardware utilisation
without inspecting logs.
